\section{Resumen y objetivos}

Este reporte final compara, con la misma partición de validación, una CNN de
referencia contra variantes que cambian un solo factor cada una. El objetivo es
responder cinco preguntas sobre la clasificación de 104 categorías de flores:

\begin{enumerate}
\item Si extraer patrones espaciales mejora la clasificación, comparando la CNN
con una red totalmente conectada (MLP) sobre las mismas imágenes.
\item Si una CNN más profunda generaliza mejor que la original.
\item Cuánto aportan la regularización (dropout, aumento de datos y decaimiento
de pesos) a la generalización.
\item Si ponderar la pérdida por clase compensa el desbalance del conjunto.
\item Qué tan sensible es el modelo a la tasa de aprendizaje y al dropout.
\end{enumerate}

Además de exactitud y F1 macro, se reportan curvas de pérdida, exactitud y F1
de entrenamiento y validación, el F1 por clase, el desbalance de los datos y
matrices de confusión. Para que la brecha de generalización sea comparable,
las métricas de entrenamiento se miden sobre imágenes sin aumento y con la red
en modo inferencia. Todas las cifras se reconstruyen con un notebook a partir de
los historiales versionados de cada corrida (Sección~\ref{sec:repo}).
